\scanpage{48}
\mspage{44}
% archive sheet number: 47

sont compatibles avec la relation \deleted{\unsure{d'ordre}} \inserted{\unsure{(}d'ensemble \deleted{\unsure{ordonné}} support}; \deleted{savoir} i.e. telles que \deleted{$\unsure{\forall x\in A\in}$} $S\in\unsure{\Phi|\varphi}$
\[
\begin{aligned}
\text{a)}\quad&S\in\Phi\Longrightarrow S\ne\varnothing_\Sigma\\
\text{b)}\quad&S,T\in\Phi,\ S\ne T\Longrightarrow S\leq CT\quad\text{i.e. }T\leq CS.
\end{aligned}
\]
\deleted{La première condition nécessaire, qui exprime la}

La condition
\[
\mathrm{D\acute{e}p}(X)\unsure{\overset{\circ}{\ll}|\overset{\circ}{\Subset}}\mathrm{D\acute{e}p}(Y)
\]
est une relation de préordre sur $\mathcal A$, que je noterai
\[
X\overset{\circ}{\ll}Y
\]
La première condition, c'est que cette relation soit une \uline{relation d'ordre}, i.e. \deleted{que}

(ii) Soient $X,Y\in\mathcal A$, tels que $\unsure{\mathrm{Dep}|\mathrm{D\acute{e}p}}\,X=\mathrm{Dep}\,Y$, alors $X=Y$.

Ceci fait, on trouve une structure $(\mathcal A,\trianglelefteq,\overset{\circ}{\ll},\mathrel{|\circ|})$, satisfaisant déjà At 1 a) b) c). Je voudrais que la relation $\ll$ qu'elle définit soit celle \deleted{de} \unsure{déduite|décrite} par $\ll$ dans $\unsure{\mathrm{Figel}\,\mathrm{ord}|\mathrm{Fig\acute{e}l}\,\mathrm{ord}|\mathrm{Fig\acute{e}l}^{\mathrm{ord}}}(\Sigma)$.

\deleted{Cela signifie que si $\mathrm{Dep}(X)$} Supposons donc \deleted{\unsure{De}}
\[
\mathrm{D\acute{e}p}(X)\ll\mathrm{Dep}(Y)
\]
donc
